% Supplementary checking notes. References use the labels of paper/paper.tex.
% This file is documentation, not an input to a proof or checker.

A cover exists: we found one with first row $M=28{,}017$ and checked it as
described in Section~\ref{sec:trust}. Its tightest inequalities are (B),
(Hc) and $Q\le(Y+\frac12)/(2(M+1))$, each satisfied with a ratio of about
$1-10^{-9}$, and (Od), with ratio about $1-2\cdot10^{-6}$.

\paragraph{What is computed, and how it is checked.}
The four premises are not proved inside Lean; they are computed.
\begin{itemize}
\item (Fa), (Fb), (Fc) are produced by two programs that share no code:
  one keeps a smallest-prime-factor table and one pointer per prime, the
  other factors by segmented sieving with a paged cursor store and
  computes $K$ by a different formula. Both run to $n=L$ and report the
  same equality cases, no violations, the same frontier thresholds, and
  identical checkpoints every $10^8$ terms: $(n,a(n),E,K,\max K)$ and a
  64-bit hash of $a(1),\dots,a(n)$. A third program, which uses neither
  pointers nor factorization, checks the defining property of every term
  up to $10^8$ (each smaller multiple of a prime factor of $a(n-1)$ appears
  earlier); swapping two terms or changing one is detected both by it and
  by the hash.
\item The cover is produced with interval arithmetic and checked by a
  separate program written from the Lean statement of
  Definition~\ref{def:cover} alone, using ball arithmetic (Arb) at 256
  bits and exact rationals and integers; every inequality is proved with
  the correct strictness. It accepts the cover and rejects corrupted ones
  (a gap, a decreased $M$, a $Q$ above its bound, $b_q=5$, $h_c=14$, a
  moved start and a short end).
\end{itemize}

\paragraph{Trusted base.}
Beyond Lean's kernel, the proof trusts: the two prefix programs, their
compilers and their recorded runs; the cover checker and the Arb library;
and the correspondence between the programs and the Lean definitions,
which was checked by reading rather than by machine: that both programs
compute \lean{EKG.actualA} by the pointer rule, that their $E$ and $K$ are
Lean's \lean{frontier} and \lean{earlyEvenCount} at the matching state
(the program's state $n$ is Lean's state $n-1$), and that the program
outputs give the premises as stated. The two prefix programs could in
principle share a misreading of the definition; the direct minimality
check covers only the first $10^8$ terms.

\paragraph{Files.}
The Lean project, the programs, the certificate and the expected outputs
accompany this paper. The Lean build takes about an hour on two cores, with up to about
$7.5$~GB of memory per Lean process. A
quick check of the premises (minutes) checks the cover and its
corruptions and replays the first $10^8$ terms; the full replay of the
first $L$ terms takes about three hours and about $72$ and $33$ GB of
memory for the two programs.


The correspondence between the statements in the paper and the formal
names is as follows.
\begin{center}
\begin{tabular}{ll}
Lemma~\ref{lem:A} & \lean{EKG3.lemmaA}\\
Corollary~\ref{cor:A} & \lean{EKG3.boundedCorollaryA}\\
Lemma~\ref{lem:C} & \lean{EKG3.lemmaC}\\
Lemma~\ref{lem:D} & \lean{EKG3.lemmaD}\\
Theorem~\ref{thm:T} & \lean{EKG3.P3.theoremT}\\
Lemma~\ref{lem:F} & \lean{EKG3.P3.lemmaF}
\end{tabular}
\end{center}


The suffix P3 in the formal names identifies the fixed prime-counting
bounds used here; it is an identifier, not a mathematical parameter.
